\documentclass[11pt,a4paper]{article}
\usepackage[T1]{fontenc}
\usepackage[utf8]{inputenc}
\usepackage{lmodern}
\usepackage[margin=25mm,headheight=14pt]{geometry}
\usepackage{amsmath,amssymb,amsthm,mathtools}
\usepackage{microtype,booktabs,longtable,array}
\usepackage{xcolor,enumitem,fancyhdr}
\usepackage{hyperref,bookmark}
\definecolor{ink}{HTML}{263B43}
\hypersetup{colorlinks=true,linkcolor=ink,citecolor=ink,urlcolor=ink,
 pdftitle={Surreal and Surcomplex Numbers in Theoretical Physics},
 pdfauthor={Research article prepared for Vladimir Reshetnikov},
 pdfsubject={Asymptotic methods, black-hole singularities, and the limits of scalar extension}}
\pagestyle{fancy}\fancyhf{}
\fancyhead[L]{\small Surreal numbers in theoretical physics}
\fancyhead[R]{\small\thepage}
\renewcommand{\headrulewidth}{0.3pt}
\setlength{\emergencystretch}{2.5em}
\setlength{\parskip}{3pt}
\setcounter{tocdepth}{2}
\numberwithin{equation}{section}
\newtheorem{theorem}{Theorem}[section]
\newtheorem{proposition}[theorem]{Proposition}
\newtheorem{lemma}[theorem]{Lemma}
\newtheorem{corollary}[theorem]{Corollary}
\theoremstyle{definition}
\newtheorem{definition}[theorem]{Definition}
\newtheorem{example}[theorem]{Example}
\theoremstyle{remark}
\newtheorem{remark}[theorem]{Remark}
\newcommand{\No}{\mathbf{No}}
\newcommand{\SC}{\No[i]}
\newcommand{\R}{\mathbb R}
\newcommand{\C}{\mathbb C}
\newcommand{\Q}{\mathbb Q}
\newcommand{\N}{\mathbb N}
\newcommand{\Z}{\mathbb Z}
\newcommand{\F}{F_{\Gamma}}
\newcommand{\K}{K_{\Gamma}}
\newcommand{\OO}{\mathcal O}
\newcommand{\mm}{\mathfrak m}
\newcommand{\Alg}{\mathcal A_{\Gamma}}
\newcommand{\eps}{\varepsilon}
\newcommand{\dd}{\mathrm d}
\newcommand{\lp}{\ell_{\mathrm P}}
\newcommand{\Krets}{\mathcal K}
\DeclareMathOperator{\st}{st}
\DeclareMathOperator{\supp}{supp}
\DeclareMathOperator{\cis}{cis}
\DeclareMathOperator{\FP}{FP}
\DeclareMathOperator{\diag}{diag}
\DeclareMathOperator{\Ric}{Ric}
\DeclareMathOperator{\PV}{PV}
\newcommand{\arxiv}[1]{\href{https://arxiv.org/abs/#1}{arXiv:#1}}
\newcommand{\doi}[1]{\href{https://doi.org/#1}{doi:#1}}
\newcommand{\repo}{\href{https://github.com/VladimirReshetnikov/Surreal/tree/39f2be6667ade51bca2b45daa47e289d69c09764}{\texttt{VladimirReshetnikov/Surreal}}}
\newcommand{\status}[1]{\textit{Status: #1.}}

\begin{document}
\begin{titlepage}
\setlength{\parindent}{0pt}
\vspace*{17mm}
{\Large\color{ink} MATHEMATICAL PHYSICS\par}
\vspace{10mm}
{\Huge\bfseries Surreal and Surcomplex Numbers in Theoretical Physics\par}
\vspace{7mm}
{\LARGE Asymptotic structure, black-hole singularities, and the limits of replacing the scalar field\par}
\vspace{13mm}
{\large Research article prepared for Vladimir Reshetnikov\par}
\vspace{3mm}
{\large September 21, 2026\par}
\vspace{12mm}
\noindent\rule{\textwidth}{0.5pt}
\begin{abstract}
Surreal numbers provide exact arithmetic for an exceptionally rich hierarchy of infinite and infinitesimal scales. Their complexification supports much of finite-dimensional complex algebra. These facts make them plausible tools for asymptotic mathematical physics, but do not by themselves make spacetime singularities regular. We distinguish scalar extension of algebraic formulas, coefficientwise non-Archimedean field theories on ordinary spacetime, and a genuinely surreal-valued spacetime theory. A detailed Schwarzschild calculation shows that replacing a small radius by a nonzero infinitesimal preserves, rather than removes, invariant curvature blow-up and destructive tidal behavior. We prove a conditional standard-part reduction theorem for Einstein's equation, identify the topological obstacles to a literal replacement of real time, and establish a finite-dimensional quantum-mechanical reduction result. Constructive applications include valuation-controlled perturbation theory, multiscale singularity diagnostics, and transseries bookkeeping, subject to support, realization, and measurement requirements. A regular-core comparison separates new dynamics from new notation. The article concludes with a concrete development and verification program for the supplied repository. This is a mathematical assessment and research proposal, not a claim to have obtained a quantum-gravity resolution of black-hole singularities.
\end{abstract}
\vfill
\noindent\small Repository snapshot: \texttt{39f2be6667ade51bca2b45daa47e289d69c09764}.\\
Established results, derivations in this article, and proposed physical hypotheses are distinguished throughout. Symbolic checks accompany the article; its proofs have not been machine-verified.
\end{titlepage}
\tableofcontents
\clearpage

\section{Assessment and the question that must actually be answered}
\label{sec:assessment}

\subsection{A qualified affirmative answer}
There is a credible mathematical role for surreal and surcomplex numbers in theoretical physics: they can encode scales, compare asymptotic contributions, preserve information that a leading-order limit discards, and support exact finite algebra. There is not, merely from their field structure, a mechanism that makes the Schwarzschild singularity traversable, restores predictability at a singular boundary, or replaces the ultraviolet dynamics absent from classical general relativity. These are different claims, and conflating them obscures both the opportunity and the difficulty.

The most productive interpretation is initially \emph{not} ``every physical number is now a surreal.'' It is: choose a controlled, set-sized non-Archimedean algebra in which physical families and their asymptotic scales can be represented, specify how differentiation and measurement work, and prove what the representation preserves. Existing work on surreal derivations, transseries embeddings, and integration gives substantial mathematical starting points, rather than only a suggestive analogy \cite{BM,ADH,CE}.

The phrase ``where usual real analysis fails'' needs refinement. Real differential geometry successfully identifies incomplete geodesics, distinguishes coordinate artifacts from invariant pathology, and estimates diverging tidal fields. The difficulty is that a particular classical spacetime may have no extension with the required regularity, or that classical dynamics ceases to be a justified physical approximation. Neither is a logical failure of real arithmetic. The role of extensions and the precise conclusions of singularity theorems are discussed carefully by Senovilla \cite{Senovilla}.

\subsection{Three different replacement proposals}
\begin{description}[style=nextline,leftmargin=0pt,labelsep=0pt]
\item[Algebraic scalar extension.]
Replace real or complex coefficients in finite algebraic expressions by elements of an ordered extension field or its complexification. Matrix identities, polynomial equations, and rational curvature formulas remain meaningful when their denominators are nonzero. This is mathematically straightforward but is not yet a theory of time evolution.
\item[Fields with non-Archimedean coefficients on real spacetime.]
Keep an ordinary spacetime domain and let fields have Hahn-series or transseries coefficients. Define spacetime derivatives coefficientwise, with explicit support conditions. This is a workable formal deformation theory and is the principal constructive framework developed below.
\item[A genuinely non-Archimedean spacetime.]
Replace coordinate domains, curves, causal structure, integration, observables, and possibly probability by new objects over a surreal field. This is a much larger proposal. A choice of field alone does not determine its topology, differential geometry, global existence theorems, or empirical interpretation.
\end{description}

These proposals can interact, but results in one do not silently become theorems in another. In particular, an expression evaluated at an infinitesimal radius is not automatically a solution on a larger spacetime containing the missing endpoint.

\subsection{What would count as success?}
A useful \emph{representation} should faithfully encode specified equations and scale relations, with computable support or remainder certificates. An \emph{analytic model} should additionally explain which formal expressions correspond to actual solutions, with which domains and boundary data. A \emph{physical singularity resolution} requires still more: an explicit dynamical law, a well-defined continuation or replacement of the singular regime, a stated regularity and measurement criterion, and physically interpretable predictions.

These standards need not all mean bounded classical curvature. A proposed quantum theory might replace spacetime geometry altogether. But then its success must be demonstrated using its own dynamics and observables, not inferred from the availability of infinite scalar values. This article evaluates mathematical possibilities, derives several useful conditional results, and proposes tests; it does not assume such a quantum theory has already been constructed.

\section{What the supplied repository actually provides}
\label{sec:repo}

The reviewed snapshot is \repo, commit \texttt{39f2be6667ad}. The inspection covered the root README, the research catalogue, the foundations at the start of the main surcomplex analysis source, and the trigonometry report's detailed description \cite{Repo,RepoAnalysis,RepoTrig}. This was a targeted mathematical review, not an independent build or an audit of every source manuscript.

The repository's most relevant contribution is its insistence on distinctions that physics cannot bypass. Its analysis separates strong Hahn summation from ordinary topological convergence. Its analytic reports distinguish fixed-domain coefficient functions, common-domain germs, individually convergent germs without a uniform domain, and entirely formal coefficient functions. Its trigonometry separates canonical finite-angle constructions from additional choices involved in global phases. These distinctions determine whether a proposed wave, integral, or perturbation is even defined.

For example, the catalogue displays
\begin{equation}
 F(z)=\sum_{n\geq 1}\frac{t^{n\eta}}{1-nz},\qquad \eta>0.
\end{equation}
Each coefficient is a holomorphic germ at zero, but its pole at $1/n$ prevents the whole family from sharing an ordinary disk of holomorphy. Treating this expression as a field on one fixed complex neighborhood would introduce a false hypothesis. In a physics application, the analogous error is to claim a uniform spacetime domain from coefficientwise local solutions \cite{Repo}.

The root README describes implemented Lean prerequisites: finite complexification and geometry, polynomial algebra and jets, a Hahn summability layer, standard part on nonnegative-order series, and compatibility between complexification and Hahn coefficients. It explicitly states that these prerequisites do not yet construct the surreal field or the normal-form bridge, and do not establish all claims in the reports. The research catalogue also labels the reports as AI-assisted, unrefereed drafts. Thus a theorem in a report and a successfully checked generic Lean declaration have different evidential status \cite{Repo}.

The proposed physics work should build on the verified algebraic interfaces where applicable, but should not cite the repository as an already formalized theory of surreal general relativity. The missing pieces include coefficient-function algebras, tensor differentiation, physical interpretation, analytic realization, and an account of evolution. Section~\ref{sec:program} makes these into concrete development targets.

\section{The scalar mathematics: powerful algebra, additional choices}
\label{sec:scalars}

\subsection{Fields, class size, and controlled workspaces}
The surreal numbers $\No$ form a real-closed ordered field in the class-sized sense. Their complexification
\begin{equation}
 \SC=\{a+ib:a,b\in\No\},\qquad i^2=-1,
\end{equation}
is algebraically closed for ordinary finite-degree polynomials. Algebraic closure is not analytic completeness. The order belongs to $\No$, not to its complexification. Conjugation and the modulus are
\begin{equation}
 \overline{a+ib}=a-ib,\qquad |a+ib|=(a^2+b^2)^{1/2}.
\end{equation}
Their multiplicative and triangle identities are algebraic consequences of real closedness \cite{RepoAnalysis,CE}.

The full class $\No$ is not a set, and neither is a putative unrestricted collection of all its field configurations. One may work in a class theory such as NBG, but every coefficient family and support used here is a \emph{set}. For calculations, choose a divisible set-sized ordered subgroup $\Gamma$ of the additive surreal group and use
\begin{equation}\label{eq:workspace}
 \F=\R((t^\Gamma)),\qquad
 \K=\C((t^\Gamma))=\F(i).
\end{equation}
Here a series $a=\sum_{\gamma\in S}a_\gamma t^\gamma$ has well-ordered support $S\subseteq\Gamma$. Such Hahn fields are compatible with Conway normal forms under the monomial embedding $t^\gamma=\omega^{-\gamma}$. A divisible $\Gamma$ gives a real-closed $\F$ and an algebraically closed $\K$; alternatively, the subsequent finite-algebra statements can simply start with any real-closed ordered field containing $\R$ \cite{RepoAnalysis}.

The notation $\omega^{-\gamma}$ here denotes Conway's monomial map. For a general surreal exponent it must not be silently identified with a Gonshor exponential expression. Exponential and logarithmic closure is extra structure, not a property of every selected Hahn workspace \cite{RepoAnalysis,vdDE,EK}.

For the rational powers in the Schwarzschild example, $\Gamma=\Q$ suffices. Arbitrary real Kasner powers can be handled in $\Gamma=\R$. Beyond-all-powers exponential scales may require an enlarged group or a specified transseries subfield. Enlarging the workspace is a mathematical operation with closure obligations, not a license to sum arbitrary families.

\subsection{Valuation, finite elements, and standard part}
For a nonzero Hahn series define
\begin{equation}
 v(a)=\min\supp(a),\qquad v(0)=+\infty.
\end{equation}
The symbol $+\infty$ in the definition of a valuation is a bookkeeping symbol, not a newly adjoined field element. We have
\begin{equation}
 v(ab)=v(a)+v(b),\qquad v(a+b)\geq\min(v(a),v(b)).
\end{equation}
A positive valuation describes an infinitesimal relative to the real coefficient field. The finite-element ring, its infinitesimal ideal, and its residue are
\begin{equation}\label{eq:standardpart}
 \OO=\{a:v(a)\geq0\},\qquad
 \mm=\{a:v(a)>0\},\qquad
 \st:\OO\longrightarrow\R\quad\text{or}\quad\C.
\end{equation}
The standard part extracts the coefficient at exponent zero and is a ring homomorphism on $\OO$. Every finite element is uniquely a real or complex constant plus an infinitesimal \cite{RepoAnalysis}.

\begin{proposition}[Two elementary obstructions]\label{prop:elementary}
No field extension makes division by zero valid. Moreover, standard part cannot extend to a unital ring homomorphism from an entire non-Archimedean field to $\R$ or $\C$ while annihilating nonzero infinitesimals.
\end{proposition}
\begin{proof}
In a field, $0b=0$ for every $b$, so $0b=1$ is impossible. For the second assertion choose a nonzero infinitesimal $\eps$. A multiplicative extension would give
$1=\st(\eps\eps^{-1})=\st(\eps)\st(\eps^{-1})=0$.
\end{proof}

Thus $1/\eps$ is meaningful and infinite, whereas $1/0$ remains undefined. Similarly, ``take the observable real part by standard part'' is a rule only for finite values, unless another measurement prescription is supplied. There is also no smallest positive infinitesimal: $\eps/2$ is a smaller one. Choosing a surreal infinitesimal does not create a minimum length.

\subsection{Dimensionful constants are not infinitesimals by being small}
The Planck length is a positive real length in conventional physics. Its small numerical value in meters does not make it infinitesimal. Statements about order should be made for dimensionless ratios, such as $r/m$, $\lp/m$, or $\lp^4\Krets$, with the reference scale specified.

If $m$ and $\lp$ are fixed nonzero real lengths and $\eps$ is a genuine infinitesimal, then $m\eps\ll\lp$, not merely ``somewhat below an astrophysical radius.'' A formal parameter representing the real-family limit $\lp/m\to0$ is useful, but identifying that parameter with a physical constant is a separate interpretive step. This distinction will matter at the effective-field-theory threshold.

\section{Topology and calculus are not inherited with the field}
\label{sec:calculus}

\subsection{Why ordinary real-time curves do not use the full fine topology}
Give $\SC$ the fine topology defined by balls of every positive surreal radius. The following fact explains the repository's caution about ordinary convergence.

\begin{proposition}[Set-discreteness of the full fine topology]\label{prop:discrete}
Every set-sized subset of $\SC$ is closed and discrete in its full fine topology. A convergent set-indexed net is eventually equal to its limit. Every continuous map from an ordinary connected real interval into this topology is constant.
\end{proposition}
\begin{proof}
For a set of distinct points, the collection of their positive pairwise distances is a set. The surreal cut construction supplies a positive number smaller than every member of this collection. Balls of a smaller radius isolate all its points. If a point lies outside the set, apply the same argument to its distances from the set to find a disjoint ball. This proves discreteness and closedness.

The range of a set-indexed net, together with its proposed limit, is again a set. A ball isolating the limit forces eventual equality. Finally, the image of an ordinary interval is a set, hence discrete; the only connected subsets of a discrete space are singletons.
\end{proof}

This is not a theorem that every notion of surreal dynamics is impossible. A class-sized surreal interval is not an ordinary real interval, and coarser topologies or coefficientwise smooth structures are different choices. Nor does Proposition~\ref{prop:discrete} apply unchanged to the \emph{intrinsic} topology of every set-sized Hahn field. For instance, $t^n\to0$ in the intrinsic valuation topology of $\R((t^\Q))$; after embedding in the full surreal field, a smaller scale such as $\omega^{-\omega}$ prevents that convergence. The two ambient topologies must not be identified.

Even in a non-Archimedean ordered field, the embedded real sequence $1/n$ does not approach zero in the full order topology: one fixed positive infinitesimal is smaller than all of its terms. A continuum model that retains real-time experiments must therefore specify what its continuity means.

\subsection{Hahn summation is neither numerical summation nor arbitrary completion}
A family of Hahn series is strongly summable when the union of its supports is well ordered and each exponent receives only finitely many contributions. In this sense
\begin{equation}\label{eq:geometric}
 \sum_{n\geq0}t^n=\frac{1}{1-t}
\end{equation}
is an exact algebraic identity. It is not asserted to be the limit of partial sums in the full fine topology. In contrast, the family of real constants $1/n!$ is not strongly Hahn-summable: all terms occupy exponent zero. The number $e$ is first defined by ordinary real analysis in the coefficient field; it is not obtained by ignoring the support rule \cite{Repo,RepoAnalysis}.

Likewise,
\begin{equation}\label{eq:factorialseries}
 \sum_{n\geq0} n!t^{n+1}
\end{equation}
is a valid Hahn series although its corresponding real power series diverges at every nonzero real argument. Exact non-Archimedean denotation and convergent real approximation are different capabilities. Section~\ref{sec:transseries} develops the physical significance of this distinction.

\subsection{Three distinct derivative constructions}
First, one can differentiate functions on a non-Archimedean field using a chosen topology and a difference-quotient definition. This does not automatically preserve familiar mean-value, compactness, or existence theorems.

Second, Berarducci and Mantova construct a compatible derivation on $\No$, with real constants, strong additivity, and compatibility with the surreal exponential; their resulting differential structure is Liouville closed \cite{BM}. Aschenbrenner, van den Dries, and van der Hoeven place it in a universal $H$-field framework with transseries and Hardy-field embeddings \cite{ADH}. Such a derivation differentiates \emph{elements interpreted as asymptotic functions}. It is not, without further work, the four commuting coordinate derivatives of a relativistic field theory. The fact that it kills real constants does not mean that a real-coordinate physical field has zero derivative: those are different mathematical objects.

Third, for a real smooth domain $U$ and a globally well-supported family, define
\begin{equation}\label{eq:coefficientderivative}
 \partial_\mu\left(\sum_\gamma a_\gamma(x)t^\gamma\right)
 =\sum_\gamma (\partial_\mu a_\gamma)(x)t^\gamma.
\end{equation}
Here the formal scales are held constant and all coefficients live on the same $U$. These derivatives commute and obey the Leibniz rule. This is the derivative used for the formal spacetime model below. It does not require nonconstant fine-continuous maps from real intervals to the full surreal class.

A change of independent variable must still obey the chain rule. If one evaluates a previously differentiated radial expression at $r=mt$, one may use the resulting algebraic value. If one instead treats $t=r/m$ as the independent radial variable, then $\partial_r=m^{-1}\partial_t$; it would be wrong to hold $t$ constant while claiming to differentiate in $r$.

\subsection{Integration and transfer: substantial results, specified scope}
Costin and Ehrlich construct surreal extensions and integrals for a large class of resurgent functions, including specified nonresonant meromorphic ODE and difference-equation solutions, and discuss foundational obstructions to much broader constructive extension principles \cite{CE}. These are important positive results. They do not amount to an unrestricted integral of every surreal-valued field over every surreal domain.

In the coefficientwise framework, integration over a fixed ordinary domain can instead be defined by integrating each coefficient, provided all coefficient integrals exist and the resulting family remains admissible. Such an operation is not the full fine-topological limit of Riemann sums. Contour integrals in the repository are explicitly of this coefficientwise kind \cite{Repo}.

Real-closed-field transfer handles first-order ordered-field statements, and stronger specified exponential structures admit stronger elementary equivalences \cite{vdDE}. Neither statement transfers the entire real analytic superstructure, its arbitrary functions, measures, and infinite-dimensional operators. A hyperreal construction with internal sets and an appropriate transfer principle is additional structure, not a consequence of having an ordered field embedding into $\No$.

\section{What a black-hole singularity means}
\label{sec:singularities}

\subsection{An endpoint is not just a very large number}
A classical Lorentzian spacetime consists of a manifold and a nondegenerate metric of Lorentzian signature, with enough differentiability for the operations under consideration. A future-directed timelike geodesic can have finite total proper time while failing to continue in that spacetime. A null geodesic can similarly have finite affine parameter. Singularity theorems concern such incompleteness under specified hypotheses; they do not generally identify a point with an infinite coordinate or assert curvature divergence in every case \cite{Senovilla}.

There are therefore several distinct questions. Does a bad coordinate chart merely need replacement? Is the spacetime extendible, and with which differentiability of the metric? Do curvature scalars or freely falling tidal components diverge? Can a chosen matter field or quantum observable continue even when classical geometry does not? Answering one does not settle the others. Artificially excising a regular point produces incompleteness without a physical curvature singularity; demanding maximality or appropriate inextendibility addresses that different issue \cite{Senovilla}.

Replacing coefficients by surreal numbers cannot supply the missing geometric definitions. A generalized theory must say what its curves, parameters, extensions, and admissible frames are. Otherwise ``the geodesic continues through an infinite value'' has no precise geometric content.

\subsection{A horizon is a useful control case}
In geometrized units $G=c=1$, let $m>0$ be the mass parameter as a length. Physical proper time is converted to length by multiplication by $c$; in conventional units $m=GM/c^2$. Schwarzschild coordinates give
\begin{equation}\label{eq:schwarzschild}
 \dd s^2=-f(r)\dd t^2+f(r)^{-1}\dd r^2+r^2\dd\Omega^2,
 \qquad f(r)=1-\frac{2m}{r}.
\end{equation}
The apparent problem at $r=2m$ is removable by ordinary real coordinates. Introduce an ingoing null coordinate $v=t+r_*$ with $\dd r_*/\dd r=f^{-1}$. Substitution gives
\begin{equation}\label{eq:ef}
 \dd s^2=-f(r)\dd v^2+2\dd v\dd r+r^2\dd\Omega^2.
\end{equation}
The radial block has determinant $-1$ and finite coefficients at $r=2m$. Thus a coordinate pole did not require any infinite number to repair.

By contrast, the invariant Kretschmann scalar of the same geometry is
\begin{equation}\label{eq:schwarzK}
 \Krets=R_{abcd}R^{abcd}=\frac{48m^2}{r^6}.
\end{equation}
Appendix~\ref{app:curvature} gives a curvature formula that derives this directly. At the horizon $\Krets=3/(4m^4)$, whereas it diverges toward zero areal radius. The vanishing of the Ricci tensor and Ricci scalar in this vacuum solution does not contradict the nonzero, divergent Weyl curvature.

Inside the horizon, $g^{ab}(\partial_a r)(\partial_b r)=f(r)<0$, so the gradient of areal radius is timelike. The future singular boundary is not analogous to an obstacle at the center of a spatial room that a traveler could circle around. Merely allowing a complex-valued radius does not preserve the real Lorentzian problem.

\subsection{The basic scalar-extension obstruction}
\begin{proposition}[Schwarzschild curvature cannot acquire a field-valued endpoint]\label{prop:noendpoint}
Let $E$ be any characteristic-zero field containing a nonzero $m$. There is no assignment of $\Krets_0\in E$ at $r=0$ that extends the identity
\begin{equation}\label{eq:invariantidentity}
 r^6\Krets=48m^2
\end{equation}
to that endpoint. In particular, adjoining surreal or surcomplex values does not extend this invariant identity through $r=0$.
\end{proposition}
\begin{proof}
At $r=0$, the left side is zero for every field element $\Krets_0$, including an element infinite relative to $\R$. The right side is nonzero.
\end{proof}

This is deliberately a precise, limited statement. It does not forbid a modified metric, a different equation, a nonclassical notion of observable, or a geometry that has no areal-radius-zero event. It disproves the specific suggestion that the same field-valued curvature identity becomes regular just because infinite elements are available. Identifying an endpoint by an added compactification symbol would instead change the algebra and would still require a metric and evolution law there.

Nieto's discussion of surreal numbers and gravitation explicitly suggests replacing the quantities in the Schwarzschild metric by surreal quantities to eliminate the singularity problem \cite{Nieto}. Proposition~\ref{prop:noendpoint} identifies the missing step in that argument. A pole at a \emph{nonzero infinitesimal} becomes an admissible infinite field element; a pole at zero has not thereby become a field-valued function, nor has an extension of spacetime been constructed. The useful intuition is the resolution of scales, not a demonstrated removal of the endpoint.

\section{A worked infalling observer: radius, curvature, and tides}
\label{sec:infall}

\subsection{Evaluating at an infinitesimal radius}
Take $m$ to be a fixed positive real length and set $r=m\eps$, with $0<\eps\ll1$. Equation~\eqref{eq:schwarzK} becomes
\begin{equation}\label{eq:epsilonK}
 \Krets(m\eps)=\frac{48}{m^4}\eps^{-6}.
\end{equation}
This value can be added, multiplied, compared, and inverted. It carries more information than a bare symbol $+\infty$: for example, it dominates $m^{-4}\eps^{-5}$ but is dominated by $m^{-4}\eps^{-7}$. It still exceeds every real curvature bound in these units. The physics of tidal exposure is not made finite by reclassifying this as a legitimate scalar.

This is already a useful positive application. The valuation of the dimensionless invariant $m^4\Krets$ records the blow-up exponent without discarding relative scale information. In a multiscale perturbation, it can distinguish leading curvature, subleading curvature, and terms that are negligible at one scale but dominant at another. Such information must still be associated with an observer and a regime of validity.

\subsection{Exact proper-time behavior for a radial geodesic}
Consider a radial timelike geodesic with conserved specific energy $E=1$. Its normalization and energy equation give
\begin{equation}
 \left(\frac{\dd r}{\dd\tau}\right)^2=E^2-f(r)=\frac{2m}{r},
 \qquad \frac{\dd r}{\dd\tau}=-\sqrt{\frac{2m}{r}}.
\end{equation}
Let $\tau_0$ be the terminal proper-time parameter in the classical solution and set $s=\tau_0-\tau>0$. Integration yields
\begin{equation}\label{eq:infall}
 s=\frac{2r^{3/2}}{3\sqrt{2m}},\qquad
 r^3=\frac{9m}{2}s^2,\qquad
 r=\left(\frac{9m}{2}\right)^{1/3}s^{2/3}.
\end{equation}
Consequently,
\begin{equation}\label{eq:properK}
 \Krets=\frac{64}{27s^4}.
\end{equation}
The cancellation of the mass in this proper-time expression is exact for this geodesic. A Hahn or surreal representation of $s$ retains the fractional power in $r$ and the fourth-order pole in curvature. It does not make the remaining proper time to the classical endpoint infinite.

\subsection{Geodesic deviation survives the substitution}
In a parallel-propagated orthonormal frame along this radial geodesic, use the convention in which the radial relative acceleration is stretching. The tidal equations are
\begin{align}
 \frac{\dd^2\xi_r}{\dd\tau^2}
 &=\frac{2m}{r^3}\xi_r=\frac{4}{9s^2}\xi_r,\label{eq:radialjacobi}\\
 \frac{\dd^2\xi_\perp}{\dd\tau^2}
 &=-\frac{m}{r^3}\xi_\perp=-\frac{2}{9s^2}\xi_\perp.
 \label{eq:transversejacobi}
\end{align}
Since $\dd^2/\dd\tau^2=\dd^2/\dd s^2$, the power-law ansatz gives
\begin{align}
 \xi_r&=A s^{-1/3}+B s^{4/3},\label{eq:radialsolutions}\\
 \xi_\perp&=C s^{1/3}+D s^{2/3}.\label{eq:transversesolutions}
\end{align}
For generic diagonal Jacobi data the radial direction stretches without bound while both transverse directions shrink; the corresponding volume scales as $s^{1/3}$ and vanishes. Exceptional choices of constants can remove the growing radial term, but do not turn these equations into a benign passage through $s=0$.

Equations~\eqref{eq:infall}--\eqref{eq:transversesolutions} are an explicit answer to the physical-process part of the question. Surreal numbers can represent the process's asymptotic scales with exact arithmetic. They preserve its destructive behavior. They do not predict what replaces that process after classical theory ceases to apply.

The asymptotic words in this paragraph refer first to ordinary real $s\downarrow0$. Substituting a nonzero infinitesimal afterward gives a Hahn-valued encoding. No convergence in the full fine surreal topology has been assumed.

\subsection{A smoother parameter is not a regular metric}
One can write $s=L\lambda^3$ with a fixed real length $L$, so that $r$ is proportional to $\lambda^2$. The fractional power in the curve disappears. But $\dd\tau/\dd\lambda=-3L\lambda^2$ vanishes at zero. This is not a nonsingular coordinate transformation at the endpoint, and the proper-time metric on the curve pulls back to
\begin{equation}
 -\dd\tau^2=-9L^2\lambda^4\dd\lambda^2,
\end{equation}
which degenerates there. Resolving a ramified formula, or blowing up an algebraic variety, does not automatically resolve a nondegenerate Lorentzian metric. Algebraic desingularization can still be a helpful organizational tool, provided this distinction is kept explicit.

\section{A precise formal model for non-Archimedean general relativity}
\label{sec:formalgr}

\subsection{The coefficient-function algebra}
Let $U\subseteq\R^4$ be a coordinate domain. Define $\Alg(U)$ to consist of series
\begin{equation}
 A(x)=\sum_{\gamma\in S}A_\gamma(x)t^\gamma,
 \qquad A_\gamma\in C^\infty(U),
\end{equation}
whose \emph{global} coefficient support $S$ is a well-ordered set in $\Gamma$. Products use Hahn convolution. This is a ring, not generally a field: smooth functions can have zeros or zero divisors. Coefficientwise differentiation is defined by \eqref{eq:coefficientderivative}. Restriction to smaller ordinary domains is allowed and does not enlarge support.

Let $\Alg^{\geq0}(U)$ and $\Alg^{>0}(U)$ denote the nonnegative- and positive-support parts. Extraction of coefficient zero gives
\begin{equation}
 \st:\Alg^{\geq0}(U)\longrightarrow C^\infty(U),
 \qquad \st\partial_\mu=\partial_\mu\st.
\end{equation}
A metric is a symmetric $4\times4$ matrix over this algebra. For an inverse and a controlled ordinary interpretation, take
\begin{equation}\label{eq:metricdeformation}
 g=g_0+h,\qquad h\in M_4(\Alg^{>0}(U)),
\end{equation}
where $g_0$ is an ordinary smooth nondegenerate Lorentzian metric on $U$. In particular, $\det g_0$ is nowhere zero. Then
\begin{equation}\label{eq:inversemetric}
 g^{-1}=\sum_{n\geq0}(-g_0^{-1}h)^n g_0^{-1}
\end{equation}
is a valid matrix Hahn sum. The positive global support and the Neumann support lemma guarantee both a well-ordered resulting support and finite coefficient contributions \cite{Repo}. This is a formal inverse, not a claim about convergence of numerical matrices.

At each ordinary point, the positive-order perturbation preserves the signature over the ordered field. One way to see this is to put $g_0$ into a real Lorentz-diagonal frame and use finite elimination: the pivots remain infinitesimal perturbations of nonzero real pivots and keep their signs. The statement is pointwise; no compactness or uniform positive lower bound on all of $U$ is being inferred.

\subsection{Levi--Civita connection and equations}
With commuting coefficient derivatives, define
\begin{equation}
 \Gamma^\alpha_{\mu\nu}
 =\frac12 g^{\alpha\beta}
 (\partial_\mu g_{\beta\nu}+\partial_\nu g_{\beta\mu}
 -\partial_\beta g_{\mu\nu}).
\end{equation}
Curvature, Ricci curvature, and the Einstein tensor are formed by their usual finite contractions and derivatives. The algebraic derivations of torsion-freeness, metric compatibility, and the Bianchi identities use only these operations and their identities. They therefore hold in this coefficientwise model. This observation is a route to formalization, not an assertion that those results are already present in the repository.

A formal Einstein equation can be written
\begin{equation}\label{eq:formalEinstein}
 G_{\mu\nu}(g)+\Lambda g_{\mu\nu}=8\pi T_{\mu\nu},
\end{equation}
with $\Lambda\in\R$ and a coefficientwise stress tensor. Conservation must either follow from the chosen matter equations or be imposed consistently with the Bianchi identity. Arbitrarily choosing coefficients of $T$ does not guarantee an admissible matter model.

\begin{theorem}[Standard-part reduction of the Einstein equation]\label{thm:reduction}
Assume \eqref{eq:metricdeformation} on one common smooth domain $U$, with $g_0$ nondegenerate everywhere, globally well-supported $h$ of strictly positive order, and $T\in M_4(\Alg^{\geq0}(U))$. Then
\begin{equation}\label{eq:Einsteinreduction}
 \st(g^{-1})=g_0^{-1},\qquad
 \st(G(g))=G(g_0).
\end{equation}
If \eqref{eq:formalEinstein} holds, its ordinary shadow satisfies
\begin{equation}
 G_{\mu\nu}(g_0)+\Lambda(g_0)_{\mu\nu}
 =8\pi\st(T_{\mu\nu}).
\end{equation}
The same coefficient-zero compatibility holds for the coordinate Riemann and Ricci tensors and their finite scalar contractions.
\end{theorem}
\begin{proof}
Only the $n=0$ term of \eqref{eq:inversemetric} has exponent zero. Standard part commutes with addition, products, real constants, and coefficient derivatives. Applying it to the displayed Christoffel formula yields the Christoffel formula for $g_0$. Apply the same operations to the finite curvature and contraction formulas. Finally reduce \eqref{eq:formalEinstein} coefficientwise.
\end{proof}

\begin{corollary}[Positive-order corrections cannot repair a regular Schwarzschild shadow]\label{cor:shadow}
A metric satisfying the hypotheses of Theorem~\ref{thm:reduction} cannot have a smooth, nondegenerate ordinary shadow through a point of zero areal radius while that shadow agrees with the Schwarzschild geometry on a punctured neighborhood with $m\ne0$.
\end{corollary}
\begin{proof}
A smooth nondegenerate ordinary metric has finite continuous curvature scalars locally. On the punctured Schwarzschild region its scalar is $48m^2/r^6$, which is unbounded as areal radius tends to zero. This contradicts such an extension.
\end{proof}

This is a conditional obstruction, not a universal theorem against non-Archimedean physics. It leaves open a degenerate or nonexistent standard-part metric, negative-order fields, singular boundary layers, new equations, nonclassical observables, or a different notion of spacetime. Any proposed evasion should identify precisely which hypothesis is being replaced and how the missing physics is supplied.

\subsection{What a formal solution does not establish}
A coefficient solution of \eqref{eq:formalEinstein} is not automatically an actual family of real solutions with that asymptotic expansion. For that conclusion one needs a realization or summation theorem, error bounds, a gauge choice, constraint propagation, and a well-posed problem in an appropriate function space. A coefficient domain that shrinks to nothing cannot support an advertised common spacetime neighborhood.

Likewise, a recursion with nonzero denominators at every formal order need not have controlled real coefficients or a convergent expansion. These issues are not reasons to reject the method; they are the distinction between formal perturbation theory and analytic existence. The successful use of surreal integration for specified function classes motivates proving such bridges in restricted relativistic models, rather than postulating them for arbitrary Einstein PDEs \cite{CE}.

\subsection{Coordinate and frame discipline}
For a finite tuple $V$, put $v_*(V)=\min_j v(V_j)$. If $A$ and $A^{-1}$ have finite entries, then
\begin{equation}\label{eq:framevaluation}
 v_*(AV)=v_*(V).
\end{equation}
Indeed, finite entries give the inequality $v_*(AV)\geq v_*(V)$, and applying $A^{-1}$ gives the reverse inequality. Tensor representations obey the same argument when the relevant frame matrices and their inverses are finite.

An infinite boost or a transformation with infinitesimal determinant need not preserve this bounded-frame criterion. Component valuations alone are therefore not unrestricted coordinate invariants. Curvature scalars, areal radius, proper time, and tidal eigenvalues in a specified observer frame are better primary diagnostics. A general theory must say which formal coordinate changes preserve its coefficient algebra and what counts as a physically comparable observer.

\section{Beyond one pole: anisotropy, weak singularities, and scale invariants}
\label{sec:other}

\subsection{Kasner geometry as a useful local benchmark}
An anisotropic vacuum benchmark is the Kasner metric
\begin{equation}\label{eq:kasner}
 \dd s_{\mathrm{geom}}^2=-\dd s^2
 +\sum_{j=1}^3(s/L_0)^{2p_j}\dd x_j^2,
 \qquad \sum_jp_j=1,\quad\sum_jp_j^2=1,
\end{equation}
where $L_0$ is a fixed length and the $x_j$ have length units. In an orthonormal frame the curvature components have coefficients proportional to $p_j(p_j-1)/s^2$ and $p_jp_k/s^2$. Squaring and counting the index symmetries gives
\begin{equation}\label{eq:kasnerK}
 \Krets=\frac{4}{s^4}
 \left[\sum_jp_j^2(p_j-1)^2+\sum_{j<k}p_j^2p_k^2\right].
\end{equation}
This formula is checked algebraically in the accompanying script. The exceptional permutations of $(1,0,0)$ give zero curvature. For $(-1/3,2/3,2/3)$ the coefficient is $64/27$, the same coefficient encountered along the Schwarzschild infall calculation.

The value of a Hahn description is not merely ``curvature is infinite.'' It simultaneously retains three anisotropic scales, their products, and the orders of tidal terms. It also makes cancellation visible: identical exponents alone do not prove a nonzero leading curvature coefficient, as the flat exceptional case demonstrates. A proposed description of a changing succession of anisotropic regimes would need additional transition dynamics and admissible supports; one Kasner expansion is not a general theorem about all black-hole interiors.

\subsection{Divergent curvature need not have the same integrated effect}
Consider the following schematic tidal magnitude, not a universal formula for black-hole interiors:
\begin{equation}\label{eq:weaktide}
 |T(s)|\asymp \frac{A}{s^2[\log(L_0/s)]^q},\qquad 0<s<s_0<L_0/e.
\end{equation}
A natural absolute-integrability diagnostic for the displacement kernel of a Jacobi equation is
\begin{equation}\label{eq:weakintegral}
 \int_0^{s_0}s|T(s)|\,\dd s
 \asymp A\int_{\log(L_0/s_0)}^\infty u^{-q}\,\dd u.
\end{equation}
The substitution $u=\log(L_0/s)$ shows that this is finite exactly when $q>1$. It is possible for curvature to diverge while this weighted tidal integral remains finite. This calculation is a diagnostic of an integral equation, not by itself a theorem giving an extension of a spacetime or its matter fields.

Logarithmic and transseries scales can distinguish these cases without collapsing them all into one infinity. However, integrability of the distortion kernel is information supplied by calculus and the equations; it does not follow merely from making $T(\eps)$ a named field element.

A concrete reason to track regularity separately is Dafermos's analysis of a spherically symmetric Einstein--Maxwell--scalar model: under its stated hypotheses the mass-inflation behavior coexists with a continuous metric extension while stronger regularity is obstructed \cite{Dafermos}. This does not settle every rotating or charged black-hole interior. It shows why a claim of resolution must specify whether it means continuous geometry, differentiable geometry, bounded curvature, or well-posed field evolution.

\subsection{Valuations classify balances, not solutions}
Suppose a dimensionless equation is written as a finite sum
\begin{equation}
 \sum_{j=1}^N c_j\prod_{k=1}^d X_k^{a_{jk}}=0,
\end{equation}
with nonzero coefficients and integer exponents where the denominators exist. If the valuations of its terms have a unique smallest value, that term cannot be canceled by the rest. Therefore any solution requires the minimum to be attained at least twice. This follows immediately from taking the leading coefficient.

This elementary valuation principle gives a concrete dominant-balance algorithm: enumerate competing leading orders, solve the resulting leading coefficient equations, and then determine corrections. It can guide asymptotic Einstein or matter equations after a gauge and ansatz have been specified. It is not a sufficiency criterion. Leading coefficients may fail to solve, constraints may conflict, logarithmic terms may be forced, or an apparently valid recursion may have no controlled real realization.

\section{The ultraviolet question is dynamical, not arithmetic}
\label{sec:eft}

\subsection{A curvature estimate for the loss of a low-energy expansion}
An effective description of gravity organizes allowed terms and their scales; changing the number field does not determine their coefficients or select a ultraviolet completion. Donoghue's formulation explains how general relativity supplies predictive low-energy quantum corrections without claiming validity at arbitrarily high curvature \cite{Donoghue}.

For the Schwarzschild example, a useful dimensionless curvature measure is
\begin{equation}
 \mathcal Q=\lp^4\Krets=\frac{48m^2\lp^4}{r^6},\qquad
 \lp=\sqrt{\frac{\hbar G}{c^3}}.
\end{equation}
The estimate $\mathcal Q\sim1$ gives
\begin{equation}\label{eq:quantumradius}
 r_{\mathrm Q}=48^{1/6}(m\lp^2)^{1/3}.
\end{equation}
This is an invariant curvature power-counting estimate, not a universal, sharply defined transition radius. Derivatives, matter scales, quantum state, and the coefficients of the effective action can introduce other relevant criteria.

For $m\gg\lp$,
\begin{equation}
 \frac{r_{\mathrm Q}}{\lp}=48^{1/6}(m/\lp)^{1/3}\gg1,
 \qquad
 \frac{r_{\mathrm Q}}{2m}
 =\frac{48^{1/6}}2(\lp/m)^{2/3}\ll1.
\end{equation}
Thus low horizon curvature for a large black hole and high interior curvature are entirely compatible. The areal radius at this estimate can be much greater than the Planck length while still being deep inside the horizon. No new scalar system is required for this separation of scales.

If $r=m\eps$ with fixed real $m,\lp>0$, then
\begin{equation}
 \mathcal Q=48(\lp/m)^4\eps^{-6}
\end{equation}
is infinite. That substitution does not keep a classical approximation under control; it moves beyond every fixed real curvature threshold. A surreal-valued computation can faithfully report that the expansion parameter is large, but cannot thereby justify neglecting higher-order terms.

\subsection{A genuinely useful two-scale calculation}
A formal large-mass family illustrates the advantage of retaining several scales. Let
\begin{equation}
 m=\lp\eps^{-a},\qquad r=m\eps^b,\qquad a,b>0.
\end{equation}
Then
\begin{equation}\label{eq:twoscale}
 \mathcal Q=48\eps^{4a-6b}.
\end{equation}
The curvature is sub-Planckian in this measure when $b<2a/3$, of order one when $b=2a/3$, and super-Planckian when $b>2a/3$. The ratio $r/m$ is infinitesimal in every case, so ``infinitesimally far inside relative to the mass scale'' alone does not determine the curvature regime when the mass is simultaneously scaled.

This does not contradict the preceding fixed-mass conclusion: it changes an explicit assumption. A valuation representation makes the competing assumptions visible and mechanically checkable. Physical application still requires choosing a finite large mass and estimating how accurately the formal family describes it.

\subsection{Hawking modes are a different issue from the central singularity}
A local mode frequency is measured relative to an observer, for example through $-k_a u^a$ in geometrized units. Large backward-traced frequencies in horizon calculations concern redshift, states, propagation, and short-distance dynamics; they are not the same invariant issue as $\Krets\to\infty$ at $r=0$.

Unruh and Sch\"utzhold derive conditions under which modified high-frequency dispersion still gives the Hawking effect to leading order, and also exhibit ways those conditions can fail \cite{Unruh}. A surcomplex symbol for an infinite phase or frequency does not choose the dispersion law, state, or boundary condition used in such a calculation. It can be a bookkeeping variable within a specified model, not a substitute for those inputs.

\section{A regular-core comparison: where new physics actually enters}
\label{sec:core}

A useful comparison is the Hayward-type static metric function \cite{Hayward}
\begin{equation}\label{eq:hayward}
 f_\ell(r)=1-\frac{2mr^2}{r^3+2m\ell^2},\qquad \ell>0,
\end{equation}
inserted into the same static spherical metric form as \eqref{eq:schwarzschild}. At large radius its correction to Schwarzschild begins at order $r^{-4}$ in $f$. At small radius,
\begin{equation}\label{eq:haywardexpansion}
 f_\ell(r)=1-\frac{r^2}{\ell^2}
 +\frac{r^5}{2m\ell^4}+O(r^8).
\end{equation}
Using the curvature formula of Appendix~\ref{app:curvature} gives
\begin{equation}\label{eq:haywardK}
 \Krets(0)=\frac{24}{\ell^4}.
\end{equation}
For a fixed positive real $\ell$ this is finite. The calculation establishes a curvature-regular central limit for this ansatz, not global geodesic completeness, stability, or a complete dynamical formation scenario.

The source of the change is evident from the mass function
\begin{equation}
 M(r)=\frac{mr^3}{r^3+2m\ell^2},\qquad f_\ell=1-\frac{2M(r)}r.
\end{equation}
If interpreted with the usual Einstein equation, its effective energy density is
\begin{equation}\label{eq:density}
 \rho(r)=\frac{M'(r)}{4\pi r^2}
 =\frac{3m^2\ell^2}{2\pi(r^3+2m\ell^2)^2},\qquad
 \rho(0)=\frac{3}{8\pi\ell^2}.
\end{equation}
This is not the vacuum Schwarzschild stress tensor. A core scale and a different geometry have entered. A physical theory must explain the effective source and its associated pressures, dynamics, perturbative stability, and connection to the exterior. Analyses of geodesically complete black-hole constructions emphasize that regularity must be assessed with global and physical conditions, rather than one finite scalar alone \cite{CompleteBH}.

Now set $\ell=m\eps$ with fixed real $m>0$ and a genuine infinitesimal $\eps$. Equation~\eqref{eq:haywardK} becomes
\begin{equation}
 \Krets(0)=24m^{-4}\eps^{-4}.
\end{equation}
The denominator at the center is algebraically nonzero, but the central curvature is not bounded by any real number. This distinction is important: ``a field-valued smooth expression exists'' and ``the curvature has a finite ordinary observable shadow'' are not the same regularity criterion.

The two real-parameter limits also illustrate nonuniformity. At every fixed $r>0$, letting the real $\ell$ tend to zero recovers Schwarzschild. Near $r^3\sim m\ell^2$, that outer approximation ceases to be uniform and the core must be resolved. A Hahn representation is well suited to organizing the outer and inner scales, but it is the altered metric or action, not surreal arithmetic, that changes the central behavior.

\section{The strongest positive use: asymptotic and transseries physics}
\label{sec:transseries}

\subsection{What is gained over writing one infinity}
A scale-resolving field can represent quantities of the form
\begin{equation}\label{eq:schematictrans}
 \sum_{\alpha,k,n}c_{\alpha k n}\,
 \eps^\alpha[\log(1/\eps)]^k\exp(-nA/\eps),
\end{equation}
provided the chosen transseries construction admits the displayed support and operations. This formula is schematic: an arbitrary triple-indexed family is not automatically admissible. In suitable classes it keeps algebraic, logarithmic, and beyond-all-orders sectors separate.

Such a language can help formulate matching conditions between asymptotic regions, compare competing terms in an equation, and retain exponentially small information that a power expansion discards. The established connections between surreal differential fields, transseries, and appropriate function extensions make this more than numerical wordplay \cite{BM,ADH,CE}. Still, much of this work can be done in a smaller Hahn or transseries field. The full proper class is usually an ambient unifying language, not a computational necessity.

\subsection{A divergent series that becomes an exact element}
Consider again
\begin{equation}
 \widehat F(z)=\sum_{n=0}^{\infty}n!z^{n+1}.
\end{equation}
For a positive infinitesimal $\eps$ it defines a Hahn element $\widehat F(\eps)$. For a nonzero real $z$, the ratio of successive term magnitudes is $(n+1)|z|$, so its ordinary series diverges. Exact denotation at $\eps$ does not give a real-valued sum at $z$.

With the Borel convention sending $n!z^{n+1}$ to $\xi^n$, the transformed series is
\begin{equation}
 \mathcal B\widehat F(\xi)=\frac1{1-\xi}.
\end{equation}
Lateral Laplace integrals bypass the pole on opposite sides:
\begin{equation}
 F_\pm(z)=\int_{0}^{\infty e^{\pm i0}}
 \frac{e^{-\xi/z}}{1-\xi}\,\dd\xi.
\end{equation}
Their difference has magnitude and phase proportional to $2\pi i e^{-1/z}$; its sign depends on the orientation convention. This follows by taking the residue of $e^{-\xi/z}/(1-\xi)$ at $\xi=1$. The two functions have the same formal power asymptotics away from the Stokes direction but differ by information invisible at every algebraic order.

A transseries must therefore include a suitable exponential sector and data selecting its coefficient. Naming the Hahn element has not selected a boundary condition, a lateral summation, or a physical state. This elementary example displays the reason that Stokes data and realization matter in resurgent physics \cite{Resurgence}.

For a genuine positive infinitesimal, $e^{-1/\eps}$ is smaller than every $\eps^N$. It is not an element of the rank-one field $\R((\eps^\Q))$ with that prescribed monomial scale: a nonzero element there has a least rational exponent. A larger workspace is needed. Extending the arithmetic and choosing how it represents an actual function are two separate steps.

\subsection{A conditional asymptotic-to-surreal pipeline}
A rigorous application can proceed through a restricted pipeline. Start with a specified family of real or complex solutions and a dimensionless parameter. Derive its admissible asymptotic expansion, with uniformity on a named ordinary domain. Embed that expansion in a differential Hahn or transseries field whose derivation represents the selected variable. Finally, prove that the observables extracted from the representation agree with the original solution up to a stated error.

There are two different possible outputs. A \emph{formal certificate} proves support admissibility and vanishing of the residual through a valuation threshold. An \emph{analytic certificate} also proves approximation or summability and quantifies a real error. The first is valuable without the second, but it must not be advertised as an actual spacetime solution.

For singular ODE reductions, existing surreal integration results suggest a feasible bridge when their hypotheses apply \cite{CE}. For the unrestricted nonlinear Einstein PDE system, that bridge would require additional work. A precise restricted model is preferable to an assertion of universal transfer.

\subsection{Series algebra can expose, rather than conceal, missing information}
A useful computational representation should store the coefficient domain, exponent group, derivative convention, and any summation prescription alongside the series. It should reject a product or substitution whose support condition fails. It should distinguish an exact Hahn equality from an asymptotic equality of functions and from equality after standard part.

For instance, a remainder of order $v(R)>N$ only bounds scale in the formal ordering. A numerical bound at a finite real parameter needs control of coefficients and a realization map. Conversely, two real solutions differing by $e^{-A/\eps}$ can have identical power-series data while remaining physically distinct in a process that amplifies that sector. Surreal richness is most useful when it prevents such information from being discarded prematurely.

\section{Surcomplex quantum mechanics: what extends and what does not}
\label{sec:quantum}

\subsection{Finite-dimensional kinematics extends algebraically}
Let $F$ be a real-closed ordered field containing $\R$, and let $K=F(i)$. For an ordinary finite dimension $d$, define
\begin{equation}
 \langle\psi,\phi\rangle=\sum_{j=1}^d\overline{\psi_j}\phi_j,
 \qquad \|\psi\|^2=\sum_{j=1}^d|\psi_j|^2.
\end{equation}
Positivity and Cauchy--Schwarz hold by the usual finite algebra. Unitary matrices preserve the norm. A normalized vector has nonnegative candidate Born weights $p_j=|\psi_j|^2$ summing to one in $F$.

\begin{proposition}[Finite Hermitian spectral theorem over a real-closed field]\label{prop:spectral}
Every Hermitian $d\times d$ matrix over $K=F(i)$ has a unitary eigenbasis and eigenvalues in $F$.
\end{proposition}
\begin{proof}
Algebraic closure of $K$ gives an eigenvalue $\lambda$ and nonzero eigenvector $v$. Since $v^\dagger v>0$,
\begin{equation}
 \lambda=\frac{v^\dagger Hv}{v^\dagger v}\in F,
\end{equation}
because the numerator is fixed by conjugation. Normalize $v$ using the positive square root of its norm. Its orthogonal complement is invariant under $H$, and the restriction is Hermitian. Induction on the ordinary finite dimension supplies an orthonormal eigenbasis.
\end{proof}

Thus finite quantum kinematics, finite spin matrices, and finite unitary algebra do not force an Archimedean scalar field. These are scalar-extension results, not new quantum predictions. An ordered extension also permits infinitesimal probabilities; a physical interpretation of those probabilities is an additional rule.

\subsection{Finite observations have an ordinary shadow}
\begin{theorem}[Reduction of finite quantum protocols]\label{thm:quantumreduction}
Suppose all dimensions and numbers of steps are ordinary finite integers. Let states be normalized over $K$, operations be unitary, and measurement probabilities be interpreted through standard part on finite elements. Then reduction gives normalized complex states, complex unitary operations, and the ordinary complex probabilities for the reduced protocol. Finite observables with finite entries have their expectation values reduced correctly. The same holds for conditioning on an outcome whose probability has nonzero standard part.
\end{theorem}
\begin{proof}
Each normalized component satisfies $|\psi_j|\leq1$, so all components are finite. Every entry of a unitary matrix is likewise finite by column normalization. Standard part commutes with conjugation, finite sums, and products. Hence
\begin{equation}
 (\st U)^\dagger\st U=I,\qquad
 \|\st\psi\|^2=1,\qquad
 \st(|\psi_j|^2)=|\st\psi_j|^2.
\end{equation}
It also commutes with the finite products defining a protocol and with $\psi^\dagger A\psi$ for finite $A$. Conditioning divides by a probability and takes a positive square root for normalization; these operations commute with reduction when the denominator's standard part is positive.
\end{proof}

The hypothesis about conditioning matters. An outcome of infinitesimal probability has zero ordinary shadow, so reduction cannot determine the conditioned state obtained by dividing by that infinitesimal. Similarly, repeated experiments with an infinite number of trials are not covered. These are precise boundaries of the equivalence, not evidence that a finite laboratory can necessarily access those operations.

\subsection{Infinite-dimensional quantum theory is not supplied by algebraic closure}
Ordinary Hilbert-space completeness, unbounded self-adjoint operators, spectral measures, strong continuity of time evolution, and countably additive probability require more structure than the field operations. A space $K^d$ with an $F$-valued norm is not automatically a Hilbert space in the usual real-normed sense. Increasing $d$ to a surreal number is not a definition of a vector space with that dimension, let alone a measure-theoretic construction.

A generalized quantum theory must explicitly choose its state spaces, topology, operator domains, and observables. A hyperfinite nonstandard model, a coefficientwise deformation of an ordinary Hilbert space, and a module over a Hahn field are different options. None inherits every useful feature of the other two merely because its scalars admit infinitesimals.

\subsection{Finite phases are canonical locally; infinite phases need structure}
For $\theta=a+\delta$ with $a\in\R$ and $\delta$ infinitesimal, a natural phase is
\begin{equation}\label{eq:finitephase}
 \cis(\theta)=e^{ia}\sum_{n\geq0}\frac{(i\delta)^n}{n!}.
\end{equation}
The ordinary coefficient $e^{ia}$ is supplied by complex analysis and the infinitesimal series is strongly Hahn-summable. The repository develops the resulting finite-angle circle theory, while Ehrlich and Kaplan supply established global exponential constructions on surcomplex structures using additional surreal structure \cite{RepoTrig,EK}. It would be incorrect to claim that global surcomplex exponentials cannot be defined.

The local rule alone, however, does not select all infinite phases. The formal series for $e^{i/\eps}$ has exponents tending without lower bound and is not an admissible Hahn Taylor substitution. One must supply a different definition or a function-level representation.

\begin{proposition}[Infinite periods of a global phase homomorphism]\label{prop:periods}
Suppose $E:(\No,+)\to\{u\in\SC:u\bar u=1\}$ is a group homomorphism extending \eqref{eq:finitephase}. Then $E$ has an infinite period.
\end{proposition}
\begin{proof}
Every unit-modulus $u$ is finite. Its complex standard part is $e^{ia}$ for some ordinary real angle $a$. Then $ue^{-ia}=1+\eta$ with infinitesimal $\eta$. The formal logarithm is defined, and conjugation with $(1+\eta)\overline{(1+\eta)}=1$ shows that this logarithm is purely imaginary. Thus $u=\cis(a+\delta)$ for a finite real-surreal angle. For any infinite $X$, choose such a finite $\theta$ representing $E(X)$. Then $E(X-\theta)=1$, and $X-\theta$ is infinite.
\end{proof}

This does not invalidate a chosen global normalization. It shows that finite-angle agreement and the homomorphism law do not reproduce the ordinary assertion that all periods are ordinary integer multiples of $2\pi$. Whether a chosen global phase construction represents a physical high-frequency limit requires a link to the phase function, its state, and its derivative, not merely to its scalar value.

\subsection{Nonuniform amplification is real mathematics, not an automatic experiment}
In units with $\hbar=1$, take $H=\eps\sigma_z$ and an evolution time $T=\eps^{-1}$. Then $HT=\sigma_z$, so its finite-phase evolution is nontrivial although $\st H=0$. Reducing $H$ before an infinite duration loses information. This simple example explains why the finite-protocol theorem does not imply equivalence under arbitrary nonuniform limits.

A proposed observable effect must nevertheless state how such a duration, coupling, postselection, or amplification is physically realized. The example demonstrates a mathematical failure of interchanging reduction and an unbounded operation. It does not establish a feasible way to observe a literal infinitesimal probability or energy shift.

\section{Quantum fields, regulators, and the meaning of an integral}
\label{sec:qft}

\subsection{Pointwise infinitesimals do not automatically produce distributions}
A familiar boundary-value identity in ordinary distribution theory is
\begin{equation}\label{eq:distribution}
 \lim_{\eta\downarrow0}\frac1{x+i\eta}
 =\PV\frac1x-i\pi\delta(x).
\end{equation}
Its content can be derived directly. Write
\begin{equation}
 \frac1{x+i\eta}=\frac{x}{x^2+\eta^2}
 -i\frac{\eta}{x^2+\eta^2}.
\end{equation}
Against a smooth compactly supported test function, the imaginary kernel integrates to $\pi$ in the limit, by $x=\eta y$ and $\int_\R(1+y^2)^{-1}\dd y=\pi$. The odd real kernel tends to the principal value after subtracting the test function's value at zero locally. This is an integral statement, not a pointwise identity of functions.

Set $\eta$ to a fixed positive infinitesimal instead. At each nonzero ordinary $x$, pointwise standard part gives $1/x$; at $x=0$ the value is infinite and standard part is unavailable. The delta contribution cannot be recovered by pointwise standard part alone. A test-function pairing, with an explicitly justified coefficientwise, asymptotic, or nonstandard integration construction, is required. Merely writing a surcomplex $i\eps$ has not supplied the distributional prescription.

This example explains a recurring issue in field theory: information can concentrate in a shrinking region whose pointwise ordinary shadow vanishes away from that region. A global integration rule must resolve the region before taking the observable shadow. The distinction is closely related to the common-domain and coherent-contour cautions in the repository, though the distribution construction itself is not asserted to be one of its proved theorems.

\subsection{Removing negative powers is not a renormalization theory}
On the ring of finite elements, standard part is multiplicative. On all Laurent series, an operation called ``finite part'' that simply extracts the exponent-zero coefficient is not. For instance,
\begin{equation}
 \FP(\eps)=\FP(\eps^{-1})=0,
 \qquad \FP(\eps\eps^{-1})=\FP(1)=1.
\end{equation}
Nor is a bare finite-part prescription invariant under a harmless-looking change of regulator coordinate. Define
\begin{equation}
 \theta=\frac{\eps}{1+a\eps},\qquad a\in\R.
\end{equation}
Then
\begin{equation}\label{eq:finitepart}
 \eps^{-1}=\theta^{-1}-a,
 \qquad \FP_{\eps}(\eps^{-1})=0,
 \qquad \FP_{\theta}(\eps^{-1})=-a.
\end{equation}
The constant assigned after subtraction depends on the prescription. Physical renormalization must specify observables, allowed counterterms, and normalization conditions; field enlargement does not fix that freedom. Effective-field-theory reasoning remains necessary for identifying which interactions and parameters have physical significance \cite{Donoghue}.

Surreal arithmetic can still help: divergent powers and logarithms can be retained symbolically, counterterm cancellation can be tested exactly, and residual dependence can be exposed. That is a potential improvement in representation and verification, not a proof that ultraviolet divergences have been physically removed.

\subsection{Actions, measures, and field configurations}
In the formal metric framework, the determinant has nonzero ordinary leading term. The square root of $-\det g$ can consequently be expanded by an admissible positive-order binomial series. On a fixed ordinary compact integration region, and with coefficient integrability specified, one may form a coefficientwise Einstein--Hilbert action and vary it with appropriate boundary conditions. Its local Euler--Lagrange expressions recover the formal equations of Section~\ref{sec:formalgr}. This is a concrete, conservative extension of classical variational algebra.

A quantum functional integral is a different task. The field space, measure or replacement for a measure, regularization, gauge treatment, and limiting procedure all need definitions. Those issues are already substantial over $\R$ and $\C$. A proper-class sum over all surreal-valued histories is not licensed by Hahn summability. Formal sums, finite-dimensional integrals, and rigorously constructed functional measures should remain explicitly different objects.

Loeb's construction is a useful comparison: it obtains an ordinary measure from specified nonstandard measure data \cite{Loeb}. Such a result relies on the surrounding nonstandard structure. An ordered-field inclusion into a surreal field does not transport all its internal sets, saturation assumptions, or measure construction. A proposal using that machinery should state it as additional infrastructure.

\section{Measurement, recoverability, and empirical novelty}
\label{sec:measurement}

\subsection{What an ordinary apparatus is supposed to report}
There are at least two interpretations of a surreal-valued observable. It can encode the asymptotic behavior of a family of ordinary observables, in which case there is a parameter-to-physical-system map and an approximation problem. Or it can be taken literally as the value in a new ontology, in which case an apparatus rule must map it to records and probabilities.

The standard-part map gives a simple apparatus rule for finite values. It is not forced by field algebra, but it is natural when infinitesimals are meant to represent deviations below every real resolution. Theorems~\ref{thm:reduction} and \ref{thm:quantumreduction} show how conservative this rule can be: regular finite data reduce to the usual theories in the stated models.

More generally, any finite expression constructed from finite inputs by addition, multiplication, conjugation, and inverses with nonzero standard part commutes with reduction. Thus a finite algebraic readout protocol cannot distinguish two inputs whose difference is infinitesimal under those conditions. This statement is mathematical, not an assertion that every imaginable physical apparatus is such a protocol.

\subsection{Ways to leave the conservative regime}
A new prediction could arise through a singular rescaling, an unbounded duration, a nonstandard observable rule, or an explicit change of dynamics. Each is a meaningful research possibility. Each changes more than the bare choice of scalar field.

For example, $x\mapsto x/\eps$ promotes an infinitesimal deviation to an order-one result. The operation is defined algebraically, but an experimental proposal must identify the physical gain, its energy or stability cost, and what fixes $\eps$. If an ordinary finite experiment with bounded couplings is claimed to see the effect, the model should derive that result without silently inserting an infinite resource.

Similarly, an exact infinitesimal probability is nonzero in the ordered field but has zero standard part. For a fixed ordinary finite number of trials, the probability of seeing such an event remains infinitesimal in a finite independent-trial model. Postselection on that event lies outside the finite-probability reduction theorem. A new probability interpretation is possible, but it has to specify how records are related to the field-valued weights.

\subsection{Invariance of predictions is a decisive test}
A physically meaningful result should not depend on arbitrary changes of presentation. At minimum, one should test changes of admissible coordinates, equivalent asymptotic parametrizations, choices of embedding into a larger Hahn workspace, and regulator prescriptions subject to the same physical normalization conditions.

This requirement is stronger than formal equality inside one chosen representation. Equation~\eqref{eq:finitepart} demonstrates how an apparently finite prediction can shift under a reparametrization when the observable rule is underspecified. Conversely, the invariant curvature identities and bounded-frame valuation statement give examples of quantities that can be controlled intrinsically.

No empirical preference for a literal surreal scalar ontology follows from the constructions in this article. Their immediate case rests on mathematical utility. An experimentally distinct theory would have to exhibit a definite finite prediction, show why ordinary effective descriptions do not already contain it, and connect its parameters to a reproducible physical procedure.

\section{Comparison with other non-Archimedean and singular methods}
\label{sec:comparison}

The choice should be driven by the operation required, not by the size of the number system. The following table summarizes the relevant distinctions; it is not a ranking of entire research programs.

\begin{longtable}{@{}>{\raggedright\arraybackslash}p{0.20\textwidth}>{\raggedright\arraybackslash}p{0.36\textwidth}>{\raggedright\arraybackslash}p{0.38\textwidth}@{}}
\toprule
Framework & Especially natural use & Additional work still required \\
\midrule
\endfirsthead
\toprule
Framework & Especially natural use & Additional work still required \\
\midrule
\endhead
Hahn fields & Exact ordered powers and support-certified arithmetic; finite algebra and valuation balances. & Function domains, derivatives, integration, and realization of formal series. \\
\addlinespace
Transseries & Logarithmic and exponential asymptotics, beyond-all-orders sectors, differential algebra. & Admissible support, summation, Stokes and boundary data, error estimates. \\
\addlinespace
Surreal numbers & A broad ambient ordered field and compatible differential or exponential structures for many scales. & Set/class discipline, a chosen analytic framework, computational representation, observables. \\
\addlinespace
Hyperreal methods & Infinitesimal encodings with a specified transfer and internal-set framework. & Appropriate internal models, standard-part operations, measure and dynamics hypotheses. \\
\addlinespace
$p$-adic models & Non-Archimedean norm geometry and explicitly constructed alternative field theories. & New geometric and physical interpretations; no inherited real ordering or Lorentzian causality. \\
\addlinespace
Distributions and generalized functions & Weak limits, singular sources, boundary values, and test-function observables. & Products or nonlinear equations may require additional structure; weak solutions need admissibility conditions. \\
\addlinespace
Regular-core or other modified dynamics & A specified change to geometry or evolution in the singular regime. & Derivation, stability, global continuation, and observable consistency. \\
\bottomrule
\end{longtable}

The surreal--transseries connections are substantive mathematical results \cite{BM,ADH,CE}. The connection to hyperreal methods is not an automatic transfer of all nonstandard analysis. Nor should $p$-adic physics be treated as surreal physics with different notation: $\Q_p$ does not carry a compatible field order, whereas order is central to the real-surreal modulus and causal-sign constructions used here. Gubser and collaborators' $p$-adic AdS/CFT model is instructive precisely because it specifies its geometry and computes correlators, rather than merely proposing a scalar substitution \cite{Padic}.

For the user's repository, a set-sized Hahn or transseries layer is the lowest-risk first implementation. It uses the existing summability and finite-algebra work while avoiding an unnecessary commitment to a full surreal-valued event manifold. A later theory can explain when embedding those calculations into $\No$ supplies a substantive additional benefit.

\section{A concrete development program for the repository}
\label{sec:program}

\subsection{Stage A: formalize a conservative local theory}
The first target should be the model of Section~\ref{sec:formalgr}, not unrestricted quantum gravity. Define a common-domain algebra of globally supported coefficient functions, prove support preservation under products and coefficient derivatives, and construct matrix inversion around a nonsingular ordinary leading matrix. Make standard part commute with these operations by theorem, not by convention.

Next define tensors, the Levi--Civita connection, curvature, and the Einstein tensor. Prove the standard-part reduction theorem and the bounded-frame valuation identity. Proposed modules might be named \texttt{Physics/CoefficientFields}, \texttt{Physics/FormalMetric}, and \texttt{Physics/StandardPartGR}; these are proposed additions, not claims about the inspected tree.

A successful output is a checked local algebraic theory with explicit assumptions. Its failure conditions should include a zero leading determinant, a non-well-ordered global support, a coefficient family lacking a common domain, and an unsupported substitution. The system should report such failures rather than assign a symbolic infinity and proceed.

\subsection{Stage B: verify exact backgrounds and observer diagnostics}
Use Schwarzschild as a mandatory control case. Verify the cancellation of the horizon-coordinate singularity, the vacuum Ricci tensor away from zero, the Kretschmann invariant, radial proper-time scaling, and the Jacobi exponents. The method must identify the horizon as removable and $r=0$ as obstructed under the unchanged invariant identity. A system that declares both regular merely because denominators become surreal symbols has failed the benchmark.

Add Kasner metrics to test multiple real exponents and vanishing leading coefficients. Add the regular-core ansatz to distinguish changed dynamics from scalar extension. Record dimensions and reference scales in the certificates. The supplied Python script checks the concrete identities used in this article, but those checks are not a Lean proof of the general framework.

\subsection{Stage C: one nontrivial matter problem, then controlled backreaction}
A defensible next experiment is a scalar field on a fixed spherically symmetric background, with a specified asymptotic region and initial or characteristic data. First derive a consistent expansion for the linear equation and compare it with ordinary asymptotic estimates. Only then attempt backreaction in the Einstein--massless-scalar system
\begin{equation}\label{eq:Einsteinscalar}
 R_{ab}=8\pi\partial_a\phi\partial_b\phi,
 \qquad \Box_g\phi=0,
\end{equation}
in geometrized units with zero cosmological constant. These equations follow from the usual minimally coupled massless scalar action and fix the normalization of $\phi$.

For example, in double-null spherical coordinates fix the convention
\begin{equation}
 \dd s^2=-\Omega^2(u,v)\dd u\dd v+r^2(u,v)\dd\Omega_2^2,
 \qquad g_{uv}=-\Omega^2/2.
\end{equation}
Then one null constraint has the form
\begin{equation}\label{eq:nullconstraint}
 \partial_u\left(\frac{\partial_u r}{\Omega^2}\right)
 =-\frac{4\pi r}{\Omega^2}(\partial_u\phi)^2,
\end{equation}
and the analogous equation holds in $v$. The mass aspect is defined geometrically by
\begin{equation}\label{eq:massaspect}
 1-\frac{2M}{r}=-\frac{4\partial_u r\partial_v r}{\Omega^2}.
\end{equation}
These formulas give explicit constraints and invariants to test, not a promise that arbitrary characteristic data admit a Hahn solution. The normalization in \eqref{eq:nullconstraint} can be checked by $R_{uu}=-2r^{-1}(r_{uu}-2\Omega_u r_u/\Omega)$ and \eqref{eq:Einsteinscalar}.

For a candidate expansion, publish the coefficient domain, scale ordering, equations solved at each order, residual valuation, and real error information when available. Check both evolution equations and constraints. A proof that only one reduced ODE is satisfied does not establish a solution of the coupled field equations.

\subsection{Stage D: prove a realization bridge}
A particularly useful result would connect a formally certified solution in a named asymptotic class to an ordinary solution with controlled remainder. In an ODE reduction this might use an existing extension or summation theorem after verifying its hypotheses. In a PDE setting it would need appropriate analytic estimates and a common domain. The selection of boundary or Stokes data should be explicit.

A substantial issue arises when an exponent varies with position. Differentiating $t^{p(x)}$ introduces $p'(x)\log t$, so the constant-exponent coefficient algebra of Section~\ref{sec:formalgr} is not automatically closed under such expressions. Either restrict the ansatz, enlarge the differential transseries class with a proof of closure, or use separate local expansions with a justified matching rule. This is the kind of concrete obstruction a research prototype should expose.

The success criterion is not merely a longer formal expansion. It is an expansion whose interpretation can be checked against a real family or against an independently defined generalized dynamics. A failure of realization can itself reveal missing sectors or inconsistent assumptions.

\subsection{Stage E: test an actual singularity-resolution hypothesis}
Only after the conservative benchmarks should one propose a modified action, constitutive law, quantum effective equation, or nonclassical observable algebra. State which classical assumptions are changed. Then test regularity in a specified sense, continuation of physical probes, preservation or replacement of constraints, stability, and ordinary-limit agreement.

For a regular-core hypothesis, specify whether the core scale is a fixed positive real, an asymptotic parameter, or a literal infinitesimal. Check the distinction made in \eqref{eq:haywardK}. For a generalized spacetime hypothesis, define its topology, curves, and notion of maximal extension before claiming completeness. For a quantum hypothesis, define the state space and its observable map before claiming unitary evolution through the singular regime.

The decisive deliverable would be a model with a definite observable consequence or a stronger mathematical existence theorem, not an expression in which the old divergent quantity has merely been assigned an infinite name. The proposed program permits both positive and negative results and provides clear criteria for recognizing each.

\section{Conclusions}
\label{sec:conclusions}

Surreal and surcomplex numbers are viable \emph{mathematical resources} for theoretical physics. Their immediate strengths are exact scale separation, finite-dimensional algebra, and the representation of controlled asymptotic expansions. Established differential and integration results make serious applications plausible, especially for restricted ODE and transseries problems. The repository's separation of algebra, summability, domains, and topology is therefore an asset rather than a technical inconvenience.

They are not a drop-in replacement for the full analytic infrastructure of $\R$ and $\C$. In the full fine topology, ordinary real-time continuous curves into $\No[i]$ are constant. Coefficientwise dynamics or a different topology can avoid this, but must be declared. Infinite-dimensional quantum theory and functional integration also require explicit constructions.

For black holes, the distinction is sharp. At a nonzero infinitesimal radius the Schwarzschild invariant has an exact infinite value; the endpoint identity still cannot hold at zero. The infall calculation preserves the fourth-order proper-time curvature pole and the stretching-and-crushing Jacobi behavior. Under regular coefficientwise assumptions, infinitesimal corrections have an ordinary Einstein shadow and cannot turn a singular Schwarzschild shadow into a smooth nondegenerate one. These are precise mathematical limits, not a dismissal of every possible non-Archimedean model.

The strongest recommendation is to use a set-sized Hahn or transseries framework as a controlled asymptotic layer, retaining ordinary spacetime and observables initially. Develop and verify its differential tensor algebra, use it to expose scales and nonuniform limits, and demand a realization theorem where an actual solution is claimed. A later singularity-resolution proposal must explain which new dynamical or geometric ingredient changes the physics. Surreal arithmetic can make that ingredient more precisely expressible; it cannot supply it by itself.

\appendix
\section{Curvature calculations and symbolic verification}
\label{app:curvature}

\subsection{Static spherical curvature}
For
\begin{equation}
 \dd s^2=-f(r)\dd t^2+f(r)^{-1}\dd r^2
 +r^2\dd\theta^2+r^2\sin^2\theta\dd\varphi^2,
\end{equation}
the independent orthonormal curvature magnitudes are $|f''|/2$, $|f'|/(2r)$ in each of the four mixed angular planes, and $|1-f|/r^2$ in the angular plane. Contraction with the curvature symmetries gives the algebraic identity
\begin{equation}\label{eq:genericK}
 \Krets=(f'')^2+4\left(\frac{f'}r\right)^2
 +4\left(\frac{1-f}{r^2}\right)^2.
\end{equation}
It is derived on a nondegenerate coordinate patch; the rational expression may then be simplified and its regular extensions tested. For $f=1-2m/r$, the three contributions are respectively $16m^2/r^6$, $16m^2/r^6$, and $16m^2/r^6$. For $f=1-r^2/\ell^2+O(r^5)$, they tend to $4/\ell^4$, $16/\ell^4$, and $4/\ell^4$.

These calculations use ordinary real differential geometry first. Substitution into a Hahn field preserves the resulting rational identities where denominators remain nonzero. No transfer theorem for arbitrary real functions is being invoked.

\subsection{What the accompanying program checks}
The archive includes \texttt{verify\_identities.py}, which uses SymPy for exact symbolic arithmetic. It independently constructs the Christoffel symbols and Riemann tensor of the general static spherical metric, contracts them to verify \eqref{eq:genericK}, checks the Schwarzschild Ricci components, and then checks the scalar, geodesic, Kasner, core, scaling, and finite-part identities used above. All 50 checks passed in the supplied run, using Python 3.13.5 and SymPy 1.14.0; the complete output is included.

The script does not verify set-theoretic foundations, global support lemmas, the standard-part reduction theorem, analytic existence, or the interpretation of quantum probabilities. It also does not construct a black-hole evolution numerically. Its role is to catch concrete algebraic mistakes in the examples, not to confer machine-checked status on the article's general proofs.

\clearpage
\section{Assumptions and claim-status ledger}
\label{app:ledger}

\begin{longtable}{@{}>{\raggedright\arraybackslash}p{0.25\textwidth}>{\raggedright\arraybackslash}p{0.43\textwidth}>{\raggedright\arraybackslash}p{0.26\textwidth}@{}}
\toprule
Claim & Essential scope or assumption & Status \\
\midrule
\endfirsthead
\toprule
Claim & Essential scope or assumption & Status \\
\midrule
\endhead
Surreal differential and integration structures & The particular structures and function classes in the cited publications. & Established literature; not unrestricted PDE transfer. \\
\addlinespace
Full fine-topology obstruction & Full $\No[i]$ topology; set-indexed nets and ordinary real intervals. & Elementary proof given here; consistent with repository analysis. \\
\addlinespace
Schwarzschild endpoint obstruction & Same invariant identity, nonzero mass, characteristic-zero field. & Exact elementary proof; not a universal no-go for new dynamics. \\
\addlinespace
Infall and tidal formulas & Radial timelike Schwarzschild geodesic with $E=1$; stated units and frame convention. & Derived explicitly; concrete identities symbolically checked. \\
\addlinespace
Einstein standard-part reduction & Common smooth domain, positive-order perturbation, nondegenerate real leading metric. & Conditional theorem proved here; not machine-verified. \\
\addlinespace
Finite quantum reduction & Ordinary finite protocol, standard-part readout, no zero-shadow conditioning denominator. & Finite algebraic theorem proved here. \\
\addlinespace
Curvature-regular core & Different metric and effective stress; fixed positive real core scale. & Local calculation, not a proof of global or physical viability. \\
\addlinespace
Hahn/transseries field models & Certified supports, chosen derivative, and explicit realization requirements. & Constructive proposed program. \\
\addlinespace
A physical resolution of black-hole singularities & New specified dynamics or geometry and an observable criterion. & Not supplied by scalar extension; not claimed in this article. \\
\bottomrule
\end{longtable}

The proofs in this article are presented as elementary deductions and a synthesis of the specified frameworks, not as claims of priority over the existing mathematical literature. The search was targeted to the supplied repository and the primary works listed below; it was not an exhaustive bibliometric survey. The document reports neither an independent Lean build of the repository nor a formal verification of its manuscript collection.

\clearpage
\begin{thebibliography}{99}
\small\raggedright
\setlength{\itemsep}{3pt}\setlength{\parskip}{0pt}
\addcontentsline{toc}{section}{References}

\bibitem{Repo}
V. Reshetnikov, \emph{Surreal}, GitHub repository, snapshot
\texttt{39f2be6667ade51bca2b45daa47e289d69c09764}.
Root README and \texttt{docs/README.md}, reviewed for this article.
\href{https://github.com/VladimirReshetnikov/Surreal/tree/39f2be6667ade51bca2b45daa47e289d69c09764}{Pinned repository snapshot}.
The research reports and the implemented Lean prerequisites have different verification status.

\bibitem{RepoAnalysis}
\emph{Surcomplex Analysis: Germs, Coherent Hahn Sections, and Infinitesimal Zero Geometry}, research draft in \texttt{Surreal},
\texttt{docs/surcomplex/analysis/article.tex}, same snapshot as \cite{Repo}.
\href{https://github.com/VladimirReshetnikov/Surreal/blob/39f2be6667ade51bca2b45daa47e289d69c09764/docs/surcomplex/analysis/article.tex}{Source and foundational conventions}.

\bibitem{RepoTrig}
\emph{Trigonometry on the Surcomplex Plane: Canonical Phases, Arbitrary-Scale Geometry, and Infinitesimal Degeneration}, research draft in \texttt{Surreal},
\texttt{docs/surcomplex/trigonometry/README.md}, same snapshot as \cite{Repo}.
\href{https://github.com/VladimirReshetnikov/Surreal/blob/39f2be6667ade51bca2b45daa47e289d69c09764/docs/surcomplex/trigonometry/README.md}{Scope and status of the trigonometry report}.

\bibitem{BM}
A. Berarducci and V. Mantova,
``Surreal numbers, derivations and transseries,''
\emph{Journal of the European Mathematical Society} \textbf{20} (2018), 339--390.
\doi{10.4171/JEMS/769}; \arxiv{1503.00315}.

\bibitem{ADH}
M. Aschenbrenner, L. van den Dries, and J. van der Hoeven,
``The surreal numbers as a universal $H$-field,''
\arxiv{1512.02267}.

\bibitem{CE}
O. Costin and P. Ehrlich,
``Integration on the Surreals,''
\arxiv{2208.14331}, revised version 5 (2024).
The article's extension and integration theorems are used only with their stated function-class restrictions.

\bibitem{vdDE}
L. van den Dries and P. Ehrlich,
``Fields of surreal numbers and exponentiation,''
\emph{Fundamenta Mathematicae} \textbf{167} (2001), 173--188.
\doi{10.4064/fm167-2-3}.

\bibitem{EK}
P. Ehrlich and E. Kaplan,
``Surreal ordered exponential fields,''
\arxiv{2002.07739}.

\bibitem{Senovilla}
J. M. M. Senovilla,
``A critical appraisal of the singularity theorems,''
\arxiv{2108.07296}; \doi{10.1098/rsta.2021.0174}.

\bibitem{Nieto}
J. A. Nieto,
``Some Mathematical and Physical Remarks on Surreal Numbers,''
\arxiv{1611.09699}; \doi{10.4236/jmp.2016.715188}.
The discussion of surreal-valued Schwarzschild quantities is a prior proposal, not an assumed singularity-resolution theorem.

\bibitem{Dafermos}
M. Dafermos,
``The interior of charged black holes and the problem of uniqueness in general relativity,''
\arxiv{gr-qc/0307013}.

\bibitem{Donoghue}
J. F. Donoghue,
``General relativity as an effective field theory: The leading quantum corrections,''
\emph{Physical Review D} \textbf{50} (1994), 3874--3888.
\doi{10.1103/PhysRevD.50.3874}; \arxiv{gr-qc/9405057}.

\bibitem{Unruh}
W. G. Unruh and R. Sch\"utzhold,
``On the Universality of the Hawking Effect,''
\arxiv{gr-qc/0408009}.

\bibitem{Hayward}
S. A. Hayward,
``Formation and evaporation of non-singular black holes,''
\emph{Physical Review Letters} \textbf{96} (2006), 031103.
\doi{10.1103/PhysRevLett.96.031103}; \arxiv{gr-qc/0506126}.

\bibitem{CompleteBH}
R. Carballo-Rubio, F. Di Filippo, S. Liberati, and M. Visser,
``Geodesically complete black holes,''
\arxiv{1911.11200}.

\bibitem{Resurgence}
I. Aniceto, G. Ba\c{s}ar, and R. Schiappa,
``A Primer on Resurgent Transseries and Their Asymptotics,''
\arxiv{1802.10441}; \doi{10.1016/j.physrep.2019.02.003}.

\bibitem{Loeb}
P. A. Loeb,
``Conversion from Nonstandard to Standard Measure Spaces and Applications in Probability Theory,''
\emph{Transactions of the American Mathematical Society} \textbf{211} (1975), 113--122.
\href{https://www.jstor.org/stable/1997222}{Journal record}.

\bibitem{Padic}
S. S. Gubser, J. Knaute, S. Parikh, A. Samberg, and P. Witaszczyk,
``$p$-adic AdS/CFT,''
\arxiv{1605.01061}; \doi{10.1007/s00220-016-2813-6}.

\end{thebibliography}
\end{document}
